\scanpage{63}
\mspage{II. 11}
% lines 1-4
nulle, \deleted{\unsure{un}} un schéma abélien sur un corps de type fini $K$ (qu'on peut $=$ supposer un corps de nbres) tel que l'action de $\operatorname{Gal}(\overline K/K)$ sur $T_\ell(A(\overline K))$ soit ess\textsuperscript{t} abélienne, \inserted{(ou seulement ess. résoluble)} soit de type CM.

% lines 5-12
1.14. On notera cependant que \deleted{les} \inserted{certains des} arguments donnés dans le présent numéro dépassent le cadre des schémas abéliens, et s'appliquent \deleted{\unsure{aussi bien}} \inserted{p.ex.} à des ``motifs'' de \deleted{degré} \inserted{poids} quelconque (une fois la théorie des motifs assise sur des bases solides). Notons seulement qu'une démonstration \unsure{de réduction|par réduction} à une situation arithmétique, ess. identique à la démonstration de 1.5, permet de prouver ceci :

% lines 13-22 and corrections between lines
\uline{Théorème} 1.15. Soient $k$ un corps sép. clos, $Y$ un schéma \deleted{connexe} \deleted{de t.f. sur $k$, \unsure{normal}} loc de t.f. sur $k$, \deleted{\unsure{qui est}} connexe et normal, $\unsure{\xi|\zeta}$ un pt géométrique de $Y$, $f:X\longrightarrow Y$ un morphisme propre \deleted{\unsure{et lisse}}, $\ell$ un nb premier $\ne\operatorname{car} k$, \deleted{\unsure{supposons}} $i$ un entier, \inserted{supposons que \deleted{le faisceau} $R^i f_*(\mathbb{Z}_\ell)$ soit un faisceau $\ell$-adique \unsure{const\textsuperscript{t}|const.} sur $Y$)} Considérons la représentation de $\pi_1(X,\xi)$ dans $E=H^i(X_\xi,\mathbb{Q}_\ell)$, \inserted{et le groupe algébrique $G\subset\operatorname{Aut}_{\mathbb{Q}_\ell}(E)$ engendré. Alors le} \deleted{Si cette représentation est \unsure{abél.} ess\unsure{.}} \deleted{\unsure{est}} \inserted{radical $R$ de $G^\circ$ est unipotent. En particulier, si} \deleted{abélienne, elle est ess. \unsure{unipotente}} \inserted{la représentation envisagée de $\pi_1(X,\xi)$ est ess. résoluble, elle est ess. unipotente; de + si on la suppose de plus semi-simple $^{(*)}$, elle} \deleted{\unsure{elle est triviale}} est essentiellement triviale. \deleted{\unsure{elle l'est}}
\margin{\unsure{tordu|tordus} (\unsure{tjs|tjrs} vrai à condition \deleted{de} \unsure{en réduisant|réduire|un réduit} $Y$)}
% slanted; note directed at the inserted hypothesis before Considérons

% lines 23-26
1.16. Dans l'énoncé 1.15, l'hypothèse de projectivité ne nous a servi que pour assurer que \deleted{\unsure{le foncteur de}} $R^i f_*(\mathbb{Z}_{\ell X})$ est un faisceau $\ell$-adique \uline{constructible} en voisinage du \unsure{point} générique $\unsure{\eta|y}$ de $Y$,

% footnote below rule
$(*)$ Ici \deleted{comme} comme dans 1.5, on conjecture \inserted{si $f$ est lisse} que l'hypothèse de $\frac12$ simplicité est automatiquement vérifiée.
% archive sheet number: 62
